\chapter{Introduction}
\chapterauthor{Florian Frohn}

The goal of this document is to serve as the canonical reference for the annual \emph{Termination and Complexity Competition (termCOMP)}.
%
Thus, it defines the various formalisms that are used at the competition, and the categories that are run for each formalism.
%
Moreover, it defines how to compute the scoring.

For each formalism, it fixes the syntax and semantics (within reason -- it does not provide formal semantics for programming languages like C or Java).
%
For each category, it defines the \emph{property} that should be analyzed.
%
In most cases, this property is ``termination'', but in some cases, more involved properties like ``almost-sure termination'' of probabilistic programs or ``runtime complexity'' are considered.
%
Moreover, for each category, it defines the valid \emph{results} and their order.

\section{Scoring}
\label{sec:scoring}

\subsection{Valid Results and Conflicts}

Each category must define two sets of valid \emph{results}: One set for proving / over-approximating the property (e.g., YES and MAYBE for termination or valid upper bounds for complexity), and one set for disproving / under-approximating the property (e.g., NO and MAYBE or valid lower bounds).
%
Furthermore, a \emph{total order} is required for each of these categories.

For example, we may have $n^0 > n^1 > \ldots$ for upper bounds, where smaller bounds are better, but $n^0 < n^1 < \ldots$ for lower bounds, where larger bounds are better.
%
Importantly, the least elements of both orders always obtain score $0$.
%
So e.g., for termination, the set of valid answers should be $\{\text{YES}, \text{MAYBE}\}$ and $\{\text{NO}, \text{MAYBE}\}$, with the orders $\text{YES} > \text{MAYBE}$ and $\text{NO} > \text{MAYBE}$, such that the result \text{MAYBE} gets score $0$.
%
Invalid outputs are interpreted as the least element of the respective set.
%
For each category, the syntax for returning valid results must be specified as well.

Finally, each category specifies which pairs of results are \emph{conflicting}.
%
E.g., for termination, YES and NO are conflicting, and for complexity, the lower bound $n^3$ conflicts with the upper bound $n^2$.

\subsection{Computing Scores}

Based on the given ordered sets of valid results $U$ and $L$, the \emph{score} of each tool is computed uniformly across categories as follows.

First, the results provided by the tools are checked for conflicts.
%
If there are conflicting results, the authors of the respective tools are asked to resolve the conflict.
%
If they fail to do so, the steering committee tries to resolve the conflict.
%
If the tool authors / the steering committee concludes that one of the answers is incorrect, then the respective tool gets penalized with the score $-10$, and the incorrect result is disregarded when computing the scores of the other tools (as described below).
%
If the conflict cannot be resolved, then the benchmark is excluded from the ranking.

Let $u_0,\ldots,u_n \in U$ with $u_i < u_{i+1}$ and $\ell_0,\ldots,\ell_m \in L$ with $\ell_i < \ell_{i+1}$ be the answers that were returned by the tools (excluding duplicates) plus the least elements of $U$ and $L$ (i.e., $u_0 = \min(U)$ and $\ell_0 = \min(L)$).
%
Then each tool that returned the result $u_i$ is awarded $\frac{i}{n}$ points, and each tool that returned the result $\ell_j$ is awarded $\frac{j}{m}$ points.\footnote{If $n$ or $m$ is $0$, then all tools are awarded $0$ points.}
%
So the maximal possible score is $2$ (e.g., for providing tight upper and lower bounds like WORST\_CASE(Omega(n), O(n))).

\subsection{Defaults for Scoring in Termination}

Unless specified otherwise, termination categories use the valid results $U = \{\text{YES}, \text{MAYBE}\}$ and $L = \{\text{NO}, \text{MAYBE}\}$, with the orders $\text{YES} > \text{MAYBE}$ and $\text{NO} > \text{MAYBE}$.
%
The only conflicting results are YES and NO.
%
So the only possible score for each benchmarks is $1$ (for answering YES or NO) or $0$ (for answering MAYBE).
%
The default syntax for results is:
\[
  \texttt{YES} \quad | \quad \texttt{NO} \quad | \quad \texttt{MAYBE}
\]

\subsection{Defaults for Scoring in Complexity Analysis}

Unless specified otherwise, complexity categories use the asymptotic classes
\[
  \mathcal{C} = \{\mathrm{Pol}_0, \mathrm{Pol}_1, \mathrm{Pol}_2, \ldots, \mathrm{Exp}, \mathrm{DExp}, \mathrm{Fin}, \omega\}
\]
Formally, for any function $f : \N \to \N \cup \{\omega\}$, its asymptotic complexity class $\iota(f) \in \mathcal{C}$ is defined by
\[
\iota(f) =
\begin{cases}
  \mathrm{Pol}_a & \text{if $a \in \N$ is the smallest number with $f(n) \in O(n^a)$,} \\
  \mathrm{Exp} & \text{if no such $a$ exists,} \\&\text{but there is a polynomial $p(n)$ with $f(n) \in O(2^{p(n)})$,} \\
  \mathrm{DExp} & \text{if no such polynomial exists,} \\&\text{ but there is a polynomial $p(n)$ with $f(n) \in O(2^{2^{p(n)}})$,} \\
  \mathrm{Fin} & \text{if no such polynomial exists,} \\&\text{ but there is no $n \in \N$ with $f(n) = \omega$,} \\
  \omega & \text{if there is an $n \in \N$ with $f(n) = \omega$.}
\end{cases}
\]
Thus, $\mathrm{Pol}_0$ denotes constant complexity, $\mathrm{Pol}_1$ linear complexity, 
$\mathrm{Pol}_a$ polynomial complexity of degree at most $a$, $\mathrm{Exp}$ single-exponential complexity, 
$\mathrm{DExp}$ double-exponential complexity, $\mathrm{Fin}$ arbitrary finite complexity 
beyond the previous classes, and $\omega$ non-termination on at least one input.

The valid scoring results are obtained from these classes by adding the fallback answer \texttt{?} for an unknown complexity.
\[
  U = L = \mathcal{C} \cup \{\texttt{?}\}
\]

For scoring complexity results, upper and lower bounds are ordered in opposite directions.
The least element is in both cases \texttt{?}.
For upper bounds, smaller classes are better, so the valid over-approximating results are ordered by
\[
  \texttt{?} < \omega < \mathrm{Fin} < \mathrm{DExp} < \mathrm{Exp} < \cdots < \mathrm{Pol}_2 < \mathrm{Pol}_1 < \mathrm{Pol}_0.
\]
Thus, a tighter upper bound receives a better score.
For lower bounds, larger classes are better, so the valid under-approximating results are ordered by
\[
  \texttt{?} < \mathrm{Pol}_0 < \mathrm{Pol}_1 < \mathrm{Pol}_2 < \cdots < \mathrm{Exp} < \mathrm{DExp} < \mathrm{Fin} < \omega.
\]
Equivalently, stronger lower bounds receive a better score.

The intended default syntax for complexity outputs is: $\texttt{WORST\_CASE(}\ell\texttt{,}u\texttt{)}$,
where $\ell, u \in \mathcal{C}$ denote the reported lower and upper asymptotic classes.

Two answers conflict if the lower bound is strictly larger than the upper bound of another answer.
For example, \texttt{WORST\_CASE(}$\mathrm{Pol}_3$\texttt{,}$\mathrm{Pol}_3$\texttt{)} 
conflicts with \texttt{WORST\_CASE(}$\mathrm{Pol}_0$\texttt{,}$\mathrm{Pol}_2$\texttt{)}.

